% Chapter6/chapter6.tex -- Conclusion and Future Work
% -----------------------------------------------------------------------------
% The conclusion is NOT a second abstract. Do not just restate what you did.
% Instead:
%   6.1 Conclusion -- reflect on whether you achieved your objectives (Section 1.2)
%       and what the significance of your findings is.
%   6.2 Future Work -- be specific and honest about what remains to be done.
%       Vague statements like "improve performance" are not useful.
%       Good future work statements identify a concrete limitation and propose
%       a specific technique to address it.

\chapter{Conclusion and Future Work}
\label{ch:conclusion}

\section{Conclusion}
\label{sec:conclusion}

This project set out to implement and optimise the inference pipeline of
TinyLlama-1.1B on the Xilinx ZCU102 FPGA platform using Vitis HLS, with
the goal of demonstrating that a billion-scale decoder-only language model
can be run on reconfigurable hardware at practical throughput and with
substantially lower power consumption than a GPU-based deployment.

Against the three objectives stated in \Cref{sec:objectives}, the outcomes are as follows.

\textit{Objective 1 -- HLS kernel implementation.} All five core transformer
operations -- RMSNorm, matrix--vector multiplication, multi-head attention,
SwiGLU, and the embedding/unembedding lookup -- were successfully implemented
as Vitis HLS kernels and verified to produce bit-accurate outputs against
a floating-point software reference on representative inputs.

\textit{Objective 2 -- INT8 quantisation.} Symmetric per-channel PTQ at
INT8 precision was applied to all weight tensors. The quantised model
retained acceptable output quality as measured by perplexity on a
held-out evaluation set, confirming that 8-bit quantisation is viable
for this model at this scale without fine-tuning.

\textit{Objective 3 -- HLS optimisation.} Application of PIPELINE, UNROLL,
and ARRAY\_PARTITION directives reduced the initiation interval of the
innermost matrix multiplication loop to II\,=\,1, achieving the theoretical
throughput ceiling for the chosen parallelism factor within the ZCU102's
DSP and BRAM budget.

\section{Recommendations for Future Work}
\label{sec:future_work}

The following directions are suggested for future work.

\textbf{Multi-layer pipelining.} The current architecture processes one
transformer layer at a time, leaving the PL idle during DMA transfers.
Overlapping the computation of layer $l$ with the DMA prefetch of layer
$l+1$'s weights using a double-buffering scheme would substantially improve
utilisation and throughput.

\textbf{Deployment on a higher-capacity platform.} Porting the design to a data-centre FPGA such as the Xilinx Alveo U200, which
provides 64\,GB of DDR4 across four independent memory channels, would remove
the memory capacity constraint entirely and allow the full floating-point weight
set to reside on-board without quantisation. The wider memory bus of the Alveo
platform would also increase the sustainable DDR bandwidth available to the
tiled matmul, reducing the load-phase stall that currently dominates per-layer
latency.
